\documentclass[20pt, a0paper, portrait]{tikzposter}

\usepackage[utf8]{inputenc}
\usepackage{lmodern} 
\usepackage{fourier} 
\usepackage{amsmath} 
\usepackage{blindtext}
\usepackage{comment}
\usepackage{enumerate}
\usepackage{xcolor}
\usepackage{boolexpr}
\usepackage{multicol}
\setlength{\columnsep}{1.4cm}

\usetikzlibrary{arrows, decorations.markings}

\definecolor{tublue}{HTML}{3571A1}
\definecolor{tulightblue}{HTML}{589DD4}
\definecolor{tuorange}{HTML}{F6C07C}

\defineblockstyle{myDefault}{
    titlewidthscale=.94, bodywidthscale=1, titleleft,
    titleoffsetx=0pt, titleoffsety=0pt, bodyoffsetx=0pt, bodyoffsety=0pt,
    bodyverticalshift=0pt, roundedcorners=0, linewidth=0.1cm,
    titleinnersep=.5cm, bodyinnersep=1cm
}{
    \begin{scope}[line width=\blocklinewidth, rounded corners=\blockroundedcorners]
       \ifBlockHasTitle %
           \draw[draw=none]%, fill=blockbodybgcolor]
               (blockbody.south west) rectangle (blocktitle.north east);
           \draw[color=tulightblue, loosely dashed]
               (blocktitle.south west) -- (blocktitle.south east);%
       \else
             \draw[draw=none]%, fill=blockbodybgcolor]
                 (blockbody.south west) rectangle (blockbody.north east);
        \fi
    \end{scope}
}

\definetitlestyle{myWave}{
    width=\paperwidth, roundedcorners=0, linewidth=0pt, innersep=1.5cm,
    titletotopverticalspace=0mm, titletoblockverticalspace=25mm,
    titlegraphictotitledistance=10pt, titletextscale=2
}{
    \coordinate (topleft) at (\titleposleft,\titlepostop);
    \coordinate (topright) at (\titleposright,\titlepostop);
    \coordinate (lefttoright) at (\titlewidth,0);
    \coordinate (head) at (0,\titlepostop-\titleposbottom);
    \coordinate (bottomright) at (\titleposright,0);
    %
    \draw[draw=none, left color=tublue, right color=tublue]%
        (topright) -- %
        (topleft) -- %
        ($(topleft) - (head)$) -- %
        ($(topright) - (head)$) -- cycle; %
    \draw[draw=none, left color=tuorange, right color=tuorange]%
        ($(topright) - (head)$ ) --  %
        ($(topleft) - (head)$) -- %
        ($(topleft) - (head)-(0,.5)$) -- %
        ($(topright) - (head)-(0,.5)$) -- cycle; %
    \draw[draw=none, left color=tulightblue, right color=tulightblue]%
        ($(topright) - (head)-(0,.5)$ ) --  %
        ($(topleft) - (head)-(0,.5)$) -- %
        ($(topleft) - (head)-(0,1.2)$) -- %
        ($(topright) - (head)-(0,1.2)$) -- cycle; %
    \draw[draw=none, left color=tublue, right color=tublue]%
        ($(topright) - (head)-(0,1.2)$ ) --  %
        ($(topleft) - (head)-(0,1.2)$) -- %
        ($(topleft) - (head)-(0,1.5)$) -- %
        ($(topright) - (head)-(0,1.5)$) -- cycle; %
    \draw[draw=none, left color=tublue, right color=tublue]%
        (bottomleft) --  %
        ($(bottomleft) + (\titlewidth, 0)$) -- %
        ($(bottomleft) + (\titlewidth, 8)$) -- %
        ($(bottomleft) + (0, 8)$) -- cycle; %
    \draw[draw=none, left color=tuorange, right color=tuorange]%
        ($(bottomleft) + (0, 8)$) -- %
        ($(bottomleft) + (\titlewidth, 8)$) -- %
        ($(bottomleft) + (\titlewidth, 8.5)$) -- %
        ($(bottomleft) + (0, 8.5)$) -- cycle; %
    \draw[draw=none, left color=tulightblue, right color=tulightblue]%
        ($(bottomleft) + (0, 8.5)$) -- %
        ($(bottomleft) + (\titlewidth, 8.5)$) -- %
        ($(bottomleft) + (\titlewidth, 9.2)$) -- %
        ($(bottomleft) + (0, 9.2)$) -- cycle; %
    \draw[draw=none, left color=tublue, right color=tulightblue]%
        ($(bottomleft) + (0, 9.2)$) -- %
        ($(bottomleft) + (\titlewidth, 9.2)$) -- %
        ($(bottomleft) + (\titlewidth, 9.5)$) -- %
        ($(bottomleft) + (0, 9.5)$) -- cycle; %
}

\definecolorpalette{sampleColorPalette} {
  \definecolor{colorOne}{named}{tublue}
  \definecolor{colorTwo}{named}{tuorange}
  \definecolor{colorThree}{named}{tulightblue}
}

\definecolorstyle{sampleColorStyle} {
  \definecolor{colorOne}{named}{blue}
  \definecolor{colorTwo}{named}{yellow}
  \definecolor{colorThree}{named}{orange}
  }{
  % Background Colors
  \colorlet{backgroundcolor}{white}
  \colorlet{framecolor}{white}
  % Title Colors
  \colorlet{titlefgcolor}{white}
  \colorlet{titlebgcolor}{tublue}
  % Block Colors
  \colorlet{blocktitlebgcolor}{white}
  \colorlet{blocktitlefgcolor}{tublue}
  \colorlet{blockbodybgcolor}{white}
  \colorlet{blockbodyfgcolor}{black}
  % Innerblock Colors
  \colorlet{innerblocktitlebgcolor}{white}
  \colorlet{innerblocktitlefgcolor}{black}
  \colorlet{innerblockbodybgcolor}{black}
  \colorlet{innerblockbodyfgcolor}{black}
  % Note colors
  \colorlet{notefgcolor}{black}
  \colorlet{notebgcolor}{colorTwo!50!white}
  \colorlet{noteframecolor}{colorTwo}
}

\definenotestyle{myNotestyle}{ targetoffsetx=0pt, targetoffsety=0pt,
  angle=45, radius=8cm, width=6cm, connection=true, rotate=0, roundedcorners=0,
  linewidth=1pt, innersep=0pt%
}{ \ifNoteHasConnection \draw[thick] (notecenter) -- (notetarget)
  node{$\bullet$}; \fi
  \draw[draw=notebgcolor,fill=notebgcolor,rotate=\noterotate] (notecenter.south
  west) rectangle (notecenter.north east); }

\definelayouttheme{sample}{
  \usecolorstyle[colorPalette=sampleColorPalette]{sampleColorStyle}
  \usebackgroundstyle{sample}
  \usetitlestyle{myWave}
  \useblockstyle{myDefault}
  \useinnerblockstyle{Minimal}
  \usenotestyle{myNotestyle}
}

\usetheme{sample}

\defineblockstyle{MyBlock}{% define a custom style for a block
    titleoffsetx=0pt, titleoffsety=0pt, bodyoffsetx=0pt, bodyoffsety=0pt,
    bodyverticalshift=0pt, roundedcorners=22, linewidth=5pt,
    titleinnersep=00mm, bodyinnersep=0mm
}{
}
% \newcommand\figureblock[3][MyBlock]{\useblockstyle{#1}\block{#2}{#3}\useblockstyle{Basic}}
\newcommand\figureblock[3][MyBlock]{\useblockstyle{#1}\block{#2}{#3}\useblockstyle{myDefault}}

% \settitle{ \vbox{ \centering \color{titlefgcolor} \Huge \textbf{\textsf{\@title}}
% \\ \@author~ \\ \@institute \\[1em] \@titlegraphic } }

% \settitle{ \vbox{ \centering \color{titlefgcolor} {\Huge \textbf{\textsf{\@title}}}
% \\ [1em] {\LARGE \textsf{\@author~}} \\ [1em] {\Large \textsf{\@institute} }} }

% \settitle{ \begin{columns} \column{0.2} \@titlegraphic \column{0.8} \vbox{ \centering \color{titlefgcolor} {\Huge \textbf{\textsf{\@title}}}
% \\ [1em] {\LARGE \textsf{\@author~}} \\ [1em] {\Large \textsf{\@institute} }} \end{columns}}

\settitle{
  \begin{minipage}{0.1\linewidth}
      \centering
      \@titlegraphic
  \end{minipage}%
  \hspace{7.5cm}
  \begin{minipage}{0.6\linewidth}
        \centering
        \color{titlefgcolor}
        {\bfseries \sffamily \Huge \@title \par}
        \vspace*{1em}
        {\LARGE \sffamily \@author \par}
        \vspace*{1em}
        {\Large \sffamily \@institute \par}
        \vspace*{1em}
        {\large \sffamily \@date}
    \end{minipage}%
  \hspace{2cm}
  \begin{minipage}{0.1\linewidth}
         \includegraphics[scale=0.85]{graphics/HSAFLogo.jpg}
    \end{minipage}%
}

\title{An enhanced soil moisture product from the ERS scatterometers}
\author{\underline{Christoph Reimer}, Isabella Pfeil and Wolfgang Wagner}
\date{ESA Living Planet Symposium, 9-13 May 2016, Prague, Czech Republic}
\titlegraphic{\includegraphics[scale=1.]{graphics/TULogo.png}}
\institute{Vienna University of Technology, Department of Geodesy and Geoinformation}
 
\begin{document}
 
  \maketitle

  \node[draw=none, rectangle, minimum width = .6cm, align=left, inner sep = 1cm,
  text=white, text width = 37cm, anchor=west] at (-42,-55)
  {\begin{enumerate}[{[1]}] \color{white} \item W. Wagner, G. Lemoine, and H.
      Rott, A method for estimating soil moisture from\\ ERS scatterometer and
      soil data, Remote Sensing of Environment, vol. 70, no. 2, pp. 191–207,
      1999. \item W. Wagner, G. Lemoine, M. Borgeaud, and H. Rott, ``A study of
      vegetation cover effects on ERS scatterometer data,'' \\ IEEE Transactions
      on Geoscience and Remote Sensing, vol. 37, no. 2II, pp. 938–948,
      1999. \item M. Vreugdenhil, W. A. Dorigo, W. Wagner, R. A. M. de Jeu, S.
      Hahn, and M. J. E. van Marle, ``Analyzing the Vegetation Parameterization
      in the TU-Wien ASCAT Soil Moisture Retrieval,'' IEEE TGRS, pp. 1–19,
      2016. \end{enumerate}};

  \node[draw=none, minimum width = 6cm, text width = 25cm, align=justify, inner
  sep = 1cm, text=white, anchor=west] at (-2,-55) {\textbf{Acknowledgments}\\ The
    authors would like to acknowledge the ESA Scatterometer Competence Center (SCIRoCCo) for their funding support. \\ \underline{http://scirocco.serco.it}};

  \node[draw=none, minimum width = 6cm, right=.5, align=right, text=white, inner sep = 1cm]
  at (26,-55) {\textbf{Christoph Reimer}\\ Vienna University of
    Technology\\ Department of Geodesy and Geoinformation\\ E-mail:
    christoph.reimer@geo.tuwien.ac.at\\ Web: http://rs.geo.tuwien.ac.at};

  \draw [color=tuorange, fill=tuorange] (-40,29.5) rectangle (-15,16.5);

  \begin{columns}

    \column{.3333} 

    \block{Introduction \phantom{p}}{Geophysical parameters derived from
      space-borne Earth Observation Systems are either assigned to discrete
      points on a fixed Earth grid (e.g. regular lat/lon grid) or located on
      orbital point nodes with a customized arrangement, often in-line with the
      instrument's measurement geometry. The driving factors of the choice and
      structure of the underlying grid are typically spatial resolution,
      instrument geometry, instrument measurement technique and application. In
      this context, the most practical solution for the choice of the sampling
      distance represents supporting the Nyquist–Shannon sampling theorem. In
      other words, the spatial sampling rate should be half of the spatial
      resolution establishing a sufficient condition to capture all the
      information from the continuous-spatial signal.\\

      In this study we propose a new global fixed Earth grid, the so-called
      Equi1 grid, and demonstrate its structure, usability and field of
      application. The grid construction follows the aforementioned condition of
      the Nyquist-Shannon sampling theorem, by applying a nearly-constant
      spatial sampling rate. The general appearance is similar to a Reduced
      Gaussian Grid, which represents a gridded coordinate system often used for
      scientific modeling on the Earth surface. The properties of the Equi1 grid
      are compared to common Earth observation grid systems (e.g. EASE-grid,
      Gaussian Grid) discussing the advantages and disadvantages. Finally, a
      real-world example will be presented based on surface soil moisture
      information derived from Metop ASCAT.}

    % \block{Grid \phantom{p}}{asdf}

    % \column{.3333} 

    % \block{Cross-over angle concept}{asdf}

    % \column{.3333} 

    % \block{Results and discussion \phantom{p}}{asdf}

    % \block{Conclusion \phantom{p}}{asdf}

  \end{columns}

\end{document}
